\documentclass[10pt,a4paper]{amsart}
\usepackage[margin=19mm]{geometry}
\usepackage{amsmath,amssymb,mathtools}
\usepackage[scr=rsfs,cal=euler]{mathalfa}
\usepackage{xfrac}
\usepackage[unicode,hidelinks,bookmarks=false]{hyperref}
\newcommand{\ie}{i.e.\ }
\newcommand{\cf}{cf.\ }
\newcommand{\resp}{respectively\ }
\newcommand{\Subsection}{\subsection}
\providecommand{\nobreakdash}{\nobreak\hbox{-}}
\setlength{\emergencystretch}{2em}
\multlinegap0pt
\usepackage{bm}
\usepackage{bbm}
\usepackage{dsfont}
\usepackage{xspace}
\usepackage{cancel}
\usepackage{mathtools}
\usepackage[shortlabels]{enumitem}
\usepackage{amsthm}
\makeatletter
\@ifundefined{theorem}{\newtheorem{theorem}{Theorem}}{}
\@ifundefined{proposition}{\newtheorem{proposition}[theorem]{Proposition}}{}
\makeatother
\makeatletter
\newcommand{\TC}{{\sf {TC}}}
\expandafter\ifx\csname myeq\endcsname\relax
\newenvironment{myeq}[1][]
{\stepcounter{figure}\begin{equation}\tag{\thefigure}{#1}}
\fi
\expandafter\ifx\csname mysubsection\endcsname\relax
\newenvironment{mysubsection}[2][]
{\begin{subsec}\begin{upshape}\begin{bfseries}{#2.}
			\end{bfseries}{#1}}
		{\end{upshape}\end{subsec}}
\fi
\newcommand{\rr}{{\mathbb {R}}}
\makeatother
\title[Parametrized Topological Complexity for a Multi-Robot System]{Parametrized Topological Complexity for a Multi-Robot System with Variable Tasks}
\author{Gopal Chandra Dutta, Amit Kumar Paul,, Subhankar Sau}
\date{Source: arXiv:2510.09323v1 (CC BY 4.0)}
\begin{document}
\maketitle

\noindent\emph{Source and licence.} Parametrized Topological Complexity for a Multi-Robot System with Variable Tasks. Version arXiv:2510.09323v1 (CC BY 4.0), distributed under the Creative Commons Attribution 4.0 International licence, as recorded in the arXiv licence metadata for this exact version and shown on the abstract page https://arxiv.org/abs/2510.09323v1. Full rights notice: licenses/arxiv-2510.09323-CC-BY-4.0.txt.\par\smallskip

\noindent\textit{Editorial scope.} Counted unit: the Proposition whose source label is \texttt{None} in the article (source0.tex lines 400--402, source label \texttt{None}), delivered together with the article's own complete original proof of it (source0.tex lines 404--424, one \texttt{proof} environment of 21 lines). No mathematics is written, repaired or re-derived: statement and proof are the article's own text line for line. Required context retained uncounted: Eq\_fibsec example (lines 289--291); Prop\_cup product example (lines 360--362); Prop\_cohom ring E23 (lines 382--390). Every definition, display and label of those lines is retained unchanged. Omitted from this package and not redistributed: 1--288, 292--359, 363--381, 391--399, 403--403, 425--1000. Nothing inside a retained excerpt is omitted or altered: every retained body is delivered exactly as the article's lines are written, so no omission receipt of any kind accompanies this package. Citations: the retained excerpts carry no citation command at all, so no bibliography entry of the article is transcribed and no reference is redistributed. Reference adaptation: none was needed and none was made; every reference inside the delivered unit resolves inside it, so no unresolved reference or citation is produced. Printed slips kept literal: no printed word, symbol, hypothesis or number of the counted statement or of its proof was altered. Layout and class: the article's own class is \texttt{amsart} with packages including amsmath, amssymb, graphicx, tikz-cd, hyperref, xy, xypic, mathtools, xspace, bm, amstext, amsfonts, euscript, amscd; the wrapper is the lane's pinned \texttt{amsart} editorial wrapper at 10pt on A4, which restates the theorem-like environments the retained text needs and adds the article's own macro definitions that the retained text needs (\texttt{\textbackslash TC}, \texttt{\textbackslash rr}), with extra wrapper packages (bm, bbm, dsfont, xspace, cancel, mathtools, enumitem). The delivered numbering is the unit's own. Assets: no retained span contains a figure environment, a graphics-inclusion command, a PDF-page inclusion, a table or a TikZ picture; no image file is copied into this package, the package declares no asset dependency and no image file is redistributed with it. Licence: version arXiv:2510.09323v1 of this article is distributed under the Creative Commons Attribution 4.0 International licence, as recorded in the arXiv licence metadata for this exact version and shown on the abstract page https://arxiv.org/abs/2510.09323v1. A full rights notice is delivered at licenses/arxiv-2510.09323-CC-BY-4.0.txt. Characters: every retained excerpt is pure ASCII, so the delivered engine, XeLaTeX with its default Latin Modern text font, reports no missing character.
\begin{equation}\label{Eq_fibsec example}
p: E=F(\rr^3, 4) \to B=F(\rr^3, 2) \text{ defined as } (o_1, o_2, z_1, z_2)\mapsto (o_1, o_2).
\end{equation}

\begin{proposition}\label{Prop_cup product example}
Let $p\colon E\to B$ be the fibration defined in \eqref{Eq_fibsec example} and $R$ be a coefficient ring. Consider the diagonal map $\Delta : E \to E^{(2, 3)}_B$ defined as $\Delta(e)= (e, e, e)$. For the homomorphism $\Delta^\ast: H^\ast(E^{(2, 3)}_B;R) \to H^\ast(E;R)$ induced by $\Delta$, if there exist cohomology classes $u_1, \ldots, u_k \in \ker[\Delta^* : H^\ast(E^{(2, 3)}_B;R) \to H^*(E; R)]$ such that $$u_1 \smile \cdots \smile u_k \neq 0 \in H^\ast(E^{(2, 3)}_B;R)$$ then $$\TC_{(2, 3)}[p: E \to B] \geq k.$$ 
\end{proposition}

\begin{proposition}\label{Prop_cohom ring E23}
The integral cohomology ring $H^*(E_B^{(2,3)})$ contains cohomology classes $w^s_{ij}$ of degree $2$, where $1\leq i < j \leq 4$ and $1\leq s \leq 3$, satisfying the following relations:
\begin{enumerate}[(a)]
\item $(w^s_{ij})^2=0$
\item $w^s_{ip} w^s_{jp}= w^s_{ij} (w^s_{jp} - w^s_{ip}) \text{ for } i<j<p$,
\item $w^s_{12} = w^{s'}_{12} \text{ for } 1\leq s \leq s' \leq 3$,
\item $w^2_{13}=w^3_{13}$ and $w^2_{23}=w^3_{23}$.
\end{enumerate}
\end{proposition}

\begin{proposition}
For the fibration in \ref{Eq_fibsec example}, we have $\TC_{(2, 3)}[p: E \to B] \geq 6.$
\end{proposition}

\begin{proof}
Recall that the diagonal map $\Delta : E \to E^{(2, 3)}_B$ is defined by $\Delta(e)= (e, e, e)$. For $1\leq s,s'\leq 3$ and $1\leq i < j \leq 4$, we have $$\Delta^*(w^s_{ij}) = (q_s\circ \Delta)^*(w_{ij}) = (q_{s'}\circ \Delta)^*(w_{ij}) = \Delta^*(w^{s'}_{ij}).$$
Thus $w^s_{ij} - w^{s'}_{ij}\in \ker(\Delta^*)$ for $1\leq s,s'\leq 3$ and $1\leq i < j \leq 4$. 

It is sufficient to show that $(w^2_{13}-w^1_{13})^2(w^2_{14}-w^1_{14})^2(w^2_{23}-w^1_{23})(w^3_{14}-w^1_{14})\neq 0.$ Note that $(w^2_{13}-w^1_{13})^2 = -2 w^2_{13} w^1_{13}$
and $(w^2_{14}-w^1_{14})^2 = -2 w^2_{14} w^1_{14}$. Thus, it is equivalent to show $w^2_{13}w^1_{13}w^2_{14}w^1_{14}(w^2_{23}-w^1_{23})(w^3_{14}-w^1_{14})\neq 0$ which is computed in the following
\begin{align*}
&~w^2_{13}w^1_{13}w^2_{14}w^1_{14}(w^2_{23}-w^1_{23})(w^3_{14}-w^1_{14}) 
%\\=& w^2_{13}w^1_{13}w^2_{14}w^1_{14}(w^2_{23}w^3_{14}-w^1_{23}w^3_{14}-w^2_{23}w^1_{14}+w^1_{23}w^1_{14})
\\= &~ w^2_{13}w^1_{13}w^2_{14}w^1_{14}(w^2_{23}w^3_{14}-w^1_{23}w^3_{14}) \quad \quad \quad (\text{since~} (w^1_{14})^2=0)
%\\=& w^2_{13}w^1_{13}w^2_{14}w^1_{14}w^2_{23}w^3_{14}-w^2_{13}w^1_{13}w^2_{14}w^1_{14}w^1_{23}w^3_{14}
\\= &~ w^2_{14}w^1_{14}w^3_{14} \{ (w^2_{13}w^2_{23})w^1_{13} - w^2_{13}(w^1_{13}w^1_{23}) \}
\\= &~ w^2_{14}w^1_{14}w^3_{14} \{ w^2_{12}(w^2_{23}-w^2_{13})w^1_{13}-w^2_{13}w^1_{12}(w^1_{23}-w^1_{13}) \} \quad (\text{using Proposition } \ref{Prop_cohom ring E23}(b))
\\= & ~ w^2_{14}w^1_{14}w^3_{14} w^2_{12}w^2_{23}w^1_{13}-
 w^2_{14}w^1_{14}w^3_{14} w^2_{12}w^2_{13}w^1_{13}-
  w^2_{14}w^1_{14}w^3_{14} w^2_{13}w^1_{12}w^1_{23}
   \\& \hspace{3cm} +w^2_{14}w^1_{14}w^3_{14} w^2_{13}w^1_{12}w^1_{13}
 \neq 0 \quad (\text{being basis elements}).
\end{align*} 
Then applying Proposition \ref{Prop_cup product example}, we get the desired result.
\end{proof}

\end{document}
